\tituloI{PLANTEAMIENTO DEL ESTUDIO}

\tituloII{Planteamiento del problema}

\parrafoII{A nivel global, la integración de la inteligencia artificial (IA) generativa en el desarrollo de software ha transformado la construcción de soluciones algorítmicas, permitiendo automatizar la escritura de código complejo. No obstante, en dominios que requieren un rigor matemático estricto, como las simulaciones físicas, investigaciones recientes sugieren que la IA puede presentar deficiencias críticas \cite{idrisov2024}. Si bien modelos de lenguaje como GPT-4 o Claude pueden producir código funcional, su capacidad para implementar con precisión leyes físicas y gestionar la estabilidad numérica es aún cuestionable \cite{niu2024efficiency}. La literatura actual indica que el código generado puede presentar errores acumulativos en los métodos de integración numérica y un rendimiento computacional subóptimo, lo que compromete su uso en software científico o de ingeniería donde la precisión es vital \cite{elnashar2024performance}.}

\parrafoII{En el contexto nacional, la adopción de herramientas como GitHub Copilot y ChatGPT en la formación de ingenieros en Perú ha crecido aceleradamente. Sin embargo, existe una brecha en la evaluación técnica de estas herramientas cuando se aplican a problemas de modelado y simulación. La dependencia de soluciones automatizadas tipo “caja negra” sin un análisis de su precisión física puede desplazar el razonamiento físico-matemático crítico \cite{indecopi2023}. Esto plantea un riesgo en la calidad del software desarrollado localmente, donde la falta de validación empírica sobre la fiabilidad del código generado para simulaciones puede derivar en sistemas con errores lógicos estructurales invisibles para el desarrollador novel \cite{siddiq2023}.}

\parrafoII{A nivel local, en la Universidad Continental sede Cusco, se observa que los estudiantes de Ingeniería de Sistemas emplean IA para resolver componentes lógicos de aplicaciones multimedia y videojuegos. No obstante, el desarrollo de motores físicos o simulaciones de movimiento se realiza generalmente sin un proceso sistemático de validación frente a modelos teóricos. Esta práctica dificulta la identificación de fallos en la fidelidad de la simulación cuando se enfrentan a escenarios de alta complejidad, generando una brecha en la formación técnica sobre la verificación y validación de software \cite{huaman2024}.}

\parrafoII{En este contexto, la unidad de análisis de la presente investigación está constituida por algoritmos de simulación física (caída libre, colisiones y fricción) generados mediante modelos de IA generativa. La problemática central radica en la ausencia de evidencia empírica que determine el efecto del uso de estas herramientas en la precisión matemática y el rendimiento computacional de dichos algoritmos, en comparación con modelos teóricos y motores de referencia como Box2D. La falta de este análisis limita la capacidad de validar la fiabilidad de la IA en contextos de ingeniería de software que exigen rigor científico.}

\tituloII{Formulación del Problema}

\tituloIII{Problema General}
\parrafoIII{¿Qué efecto tiene el uso de herramientas de inteligencia artificial generativa en la precisión matemática y el rendimiento computacional de simulaciones físicas básicas?}

\tituloIII{Problemas Específicos}
\begin{listaPE}
    \item ¿Qué efecto tiene la generación de código mediante IA en el error cuadrático medio (MSE) de trayectorias físicas?
    \item ¿Cómo influye la complejidad del escenario físico generado por IA en el tiempo de ejecución y consumo de memoria?
    \item ¿Qué nivel de intervención técnica (correcciones) requiere el código generado por IA para validar su estabilidad numérica?
\end{listaPE}

\tituloII{Objetivos}

\tituloIII{Objetivo General}
\parrafoIII{Evaluar el efecto del uso de herramientas de inteligencia artificial generativa en la precisión matemática y el rendimiento computacional de simulaciones físicas básicas.}

\tituloIII{Objetivos Específicos}
\begin{listaOE}
    \item Medir el efecto de la generación de código por IA en el error cuadrático medio de trayectorias físicas respecto al modelo teórico.
    \item Determinar el impacto de la complejidad del escenario físico en el rendimiento computacional del código generado por IA.
    \item Validar la estabilidad numérica del código generado por IA mediante la cuantificación de las correcciones técnicas requeridas.
\end{listaOE}

\tituloII{Justificación de la investigación}

\tituloIII{Justificación Teórica}
\parrafoIII{La investigación se sustenta en los principios de la física computacional y la verificación de software. Busca contrastar la capacidad de razonamiento lógico-matemático de los modelos de IA frente a las leyes de la mecánica clásica. Aportará datos sobre la fidelidad de los métodos de integración numérica (como Euler o Verlet) implementados por la IA, contribuyendo al conocimiento sobre los límites de la IA en dominios científicos.}

\tituloIII{Justificación Práctica}
\parrafoIII{Proporciona criterios técnicos para desarrolladores que integran IA en la creación de videojuegos o simuladores médicos. Permitirá identificar si el código generado es apto para entornos de producción o si introduce riesgos de inestabilidad que degraden la experiencia del usuario o la precisión del sistema.}

\tituloIII{Justificación Metodológica}
\parrafoIII{Propone un protocolo de evaluación experimental basado en el error absoluto y el benchmarking. Establece una métrica de comparación entre la salida de un modelo de lenguaje y un motor físico estándar (Box2D), creando un flujo de trabajo replicable para futuras evaluaciones de nuevos modelos de IA.}

\tituloII{Delimitación de la investigación}

\tituloIII{Delimitación Espacial}
\parrafoIII{El estudio se desarrollará en un entorno controlado bajo Arch Linux en la Universidad Continental, sede Cusco, utilizando hardware de arquitectura x86\_64 para eliminar ruidos de virtualización.}

\tituloIII{Delimitación Temporal}
\parrafoIII{La fase experimental y recolección de métricas se ejecutará entre julio y septiembre de 2026.}

\tituloII{Hipótesis y Variables}

\tituloIII{Hipótesis}
\parrafoIII{\textbf{Hipótesis General:} El uso de herramientas de IA generativa produce cambios significativos (errores acumulativos e ineficiencias) en la precisión y rendimiento de las simulaciones físicas.}
\parrafoIII{\textbf{Hipótesis Específicas:}}
\begin{listaHE}
    \item El código generado por IA presenta un error cuadrático medio significativamente mayor al modelo matemático estándar.
    \item El incremento en la complejidad del escenario físico generado por IA produce un aumento no optimizado en el tiempo de ejecución.
    \item El código de simulación generado por IA requiere intervenciones técnicas manuales para corregir inestabilidades numéricas.
\end{listaHE}

\tituloIII{Definición de Variables}
\parrafoIII{\textbf{Variable Independiente:} Herramientas de IA generativa (uso y manipulación de prompts).}
\parrafoIII{\textbf{Variables Dependientes:}}
\begin{listaVD}
    \item Precisión matemática: medida a través del error cuadrático medio (MSE) de las variables de estado (posición y velocidad) frente al modelo teórico.
    \item Rendimiento computacional: medido a través del tiempo de ejecución (usando Hyperfine/GNU time) y el consumo de memoria RAM.
\end{listaVD}